%% 第31回編集委員会（物理第5巻）：組 原W7b
%% p5-00.tex  原-0　序章　スマホと原子炉、どちらも原子の物理で動く？（★1〜2。3段にしない：N-1）
%% 中身は計画書 tools/phys5-plan.md の B「原-0」。本文の章（原-1〜15）を読んで書いた、この巻の地図。第4巻の p4-00.tex・第3巻の p3-00.tex と同じ形。
%% 序章なので、じぶんでとこう・章末問題はない（例題は7段の見本3つ）。「この章で導く中学・高校の公式」の欄は置かない
%%   （A2 の「原子はとても小さい」は原-1 ⑦が持ち主、A8 の「エネルギーを得ている」は原-15 ①・原-16 ①。第19回 6-1-2・K-8）。
%%   はてな？の答え（1円玉の原子の数と大きさ）は①・例題1。①に、質量→原子の数→1個の体積→1辺の流れ図（phflow）を1組だけ置いた（計画書 B 原-0 の3）。
%% 記号カードは置かない（アボガドロ定数 N_A は原-1 ⑦、プランク定数 h は原-2、電子ボルトは電-2 ⑥。U-6・X-2）。
%% 記号（この章）：質量 m〔g〕、1 mol の質量（モル質量）M〔g/mol〕、密度 ρ〔g/cm³〕、アボガドロ定数 N_A、原子の数 N、1個の体積 v、1辺 a、
%%   エネルギー E〔eV・J〕、長さ L。この章では a を「1個の場所の立方体の1辺」の意味だけで使う（例題3 のケイ素の格子定数は文字にしない）。
%% 例題の並び：計画書の例題1（1円玉：はてな？）・2（燃える と 分かれる）・3（チップの線にケイ素の原子）をそのまま①②③に置いた。
%% 数値はすべて Python で計算した：checks/物理第5巻/原-0/calc_p5_00.py（出力 calc_p5_00_out.txt）。常用漢字の確かめ：kanji_p5_00.py。
%% 本文の章とそろえた数値（同じ題材は、持ち主の章の値。計画書 C-4 の「同じ題材の数をそろえる」）：
%%   ・1円玉 2.2×10²² 個・1辺 約0.26 nm（原-0 が持ち主）。原子や小さな分子の大きさ 0.1〜0.3 nm・水の分子 0.31 nm：原-1 ⑦・例題8。
%%   ・炭素1個が燃える 約4.1 eV（燃焼熱 393.5 kJ/mol）、核分裂 200 MeV、約5千万倍、1 kg で約250万倍：原-15 ①・例題1（1 g で 8.2×10¹⁰ J）、熱-16。
%%   ・300 K の k_BT＝0.0259 eV、hc＝1240 eV·nm、可視光の光子 1.6〜3.3 eV：原-2 の桁の感覚。水素の電離 13.6 eV：原-4 ⑥。
%%   ・X線写真の X線 60 keV：原-13 ③（はてな？の答え）。セシウム137 のγ線 0.662 MeV：原-13・原-14。
%%   ・原子核の半径 約 1.2A^{1/3} fm、金の原子核は原子の約2万分の1：原-1 ⑤。電子の大きさは測れる限り0：原-7 ⑥。
%%   ・ケイ素の格子定数 0.5431020511 nm（CODATA 2022）、密度 2.329 g/cm³：原-8 ③。チップの細い所が原子数十個ほど・MOSFET が100億個をこえる：原-10 ②⑦⑫。
%%   ・光の顕微鏡の限界 約0.2 μm：波-13・原-1 のやってみようの答え・原-3 ⑥。可視光 約380〜780 nm：波-13。
%%   ・数学 2-7 のはてな？「紙を42回折ると、月まで届く？」：やってみよう（半分に切っていく）の向きを逆にした題材。
%%   ・各章のはてな？の答えの節は、書き上がった原稿で確かめた（原-1 ⑤、原-2 ⑪、原-3 ⑥、原-4 ③、原-5 ④、原-6 ④、原-7 ⑨、原-8 ⑨、
%%     原-9 ⑨、原-10 ⑦、原-11 ⑪、原-12 ⑦、原-13 ③、原-14 ②、原-15 ⑨）。原-16 はまだ書かれていないので、計画書 B の②③で指した。
%%   ・各章の「まだわかっていないこと」の題は、書き上がった原稿の phmada の題で確かめた。
%%   ・④の年表の年は、各章の本文に書かれた年（とその章の出典）にそろえた：1895 原-13 ①、1896 原-12 ①、1897 原-1 ②、1900 原-2 ①、
%%     1905 原-2 ③・原-11 のコラム、1911 原-1 ⑤、1913 原-4 ⑤、1924 原-3 ①、1926 原-5 ⑥、1932 原-1 ⑥・原-11 ③、1938 原-15 ①、
%%     1942 原-15 のコラム、1947 原-10 のコラム。③の「身近な例では」の表も、各章の本文にある題材だけ（原-2 ⑫・原-7 ⑥・原-11 ⑦・原-13 ④）。
%% 出典（手引き 2章の16：書くときに確かめた。11月11日）
%%   ・1円玉：通貨の単位及び貨幣の発行等に関する法律施行令（昭和63年政令第50号）別表第一「一円の貨幣　素材 アルミニウム　品位 純アルミニウム
%%       量目 一グラム　形式 直径 二十ミリメートル」。e-Gov 法令 API https://laws.e-gov.go.jp/api/1/lawdata/363CO0000000050 で本文を読んだ。
%%   ・アルミニウムの標準原子量 26.9815385：NIST "Atomic Weights and Isotopic Compositions"（https://physics.nist.gov/cgi-bin/Compositions/stand_alone.pl?ele=Al）。
%%     密度 2.699 g/cm³：NIST "X-Ray Mass Attenuation Coefficients" 表1（https://physics.nist.gov/PhysRefData/XrayMassCoef/tab1.html）。どちらも curl で読んだ。
%%   ・定数：CODATA 2022（https://physics.nist.gov/cuu/Constants/Table/allascii.txt を curl で読んだ）：N_A＝6.02214076×10²³ /mol・e＝1.602176634×10⁻¹⁹ C
%%       （定義値）、ケイ素の格子定数 5.431020511×10⁻¹⁰ m。
%%   ・炭素の燃焼熱：CO₂（気体）の標準生成エンタルピー −393.51±0.13 kJ/mol（NIST Chemistry WebBook、CODATA Review の値。
%%       https://webbook.nist.gov/cgi/cbook.cgi?ID=C124389&Mask=1 を curl で読んだ）。本文は熱-16・原-15 と同じ 393.5 kJ/mol。
%%   ・月までの平均の距離 384400 km：NASA Moon Fact Sheet（第1巻 力-6 の頭のコメントと同じ出典・同じ値）。
%%   ・電子の大きさ：電子の磁石の強さ（磁気モーメント）を 0.13 ppt ではかり、点の電子として計算した標準模型の予言と1兆分の1のけたまで
%%       くらべた：X. Fan ほか, "Measurement of the Electron Magnetic Moment", Phys. Rev. Lett. 130, 071801 (2023)、arXiv:2209.13084 の要旨を読んだ。
%%       本文は「大きさが見つかっていない（原-7 ⑥）」「点として計算した値と1兆分の1のけたまで合う」とだけ書いた。
%%   ・ウラン235 の1回の分裂のエネルギー 約200 MeV：原-15 ①（Algora ら 2021 の内訳 202.8 MeV）と同じ。
%%   ・〔要確認〕なし（Wikipedia は出典にしていない）。チップの細い所の幅 10 nm は例題のための仮の値（本文に「〜とします」と書いた）。
\makeatletter
% 会話の1行  \txatalk{y}{sora|mina|delta}{名前}{せりふ}（p4-00 の \sxatalk と同じ形）
\newcommand\txatalk[4]{%
  \bkChar{#2}{(7,#1-6.5)}{12mm}
  \node[rectangle callout, callout absolute pointer={(15.5,#1)}, anchor=west, draw=bkTheme!80!black, line width=0.8pt,
    fill=white, rounded corners=2.5mm, inner xsep=3mm, inner ysep=1.8mm, text width=103mm, align=left,
    font=\small] at (21,#1) {#4};
  \node[anchor=north, font=\sffamily\footnotesize, text=bkGray, inner sep=0pt] at (7,#1-7.5) {#3};}
% 原子の玉  \txaatom{(中心)}{半径}{色}
\newcommand\txaatom[3]{\shade[ball color=#3!45] #1 circle (#2);}
% 立方体（ななめから見た形）  \txacube{(左下手前)}{1辺}{色}
\newcommand\txacube[3]{\begin{scope}[shift={#1}]
  \filldraw[fill=#3!18, draw=#3!70!black, line width=0.5pt] (0,0) rectangle (#2,#2);
  \filldraw[fill=#3!32, draw=#3!70!black, line width=0.5pt] (0,#2) -- ({#2*0.4},{#2*1.3}) -- ({#2*1.4},{#2*1.3}) -- (#2,#2) -- cycle;
  \filldraw[fill=#3!8, draw=#3!70!black, line width=0.5pt] (#2,0) -- ({#2*1.4},{#2*0.3}) -- ({#2*1.4},{#2*1.3}) -- (#2,#2) -- cycle;
\end{scope}}
\makeatother
\tikzset{txa lab/.style={font=\sffamily\footnotesize, inner sep=1pt, text=bkInk},
  txa arr/.style={-{Stealth[length=2.6mm, width=2.2mm]}, line width=1.1pt},
  txa box/.style={bk node, minimum width=28mm, text width=26.5mm, minimum height=18mm, font=\sffamily\footnotesize, align=center, inner xsep=1pt, inner ysep=1.6mm},
  txa dim/.style={bkGray!80!black, line width=0.5pt, {Stealth[length=1.8mm]}-{Stealth[length=1.8mm]}}}

\phchapter{0}{スマホと原子炉、\\どちらも原子の物理で動く？}{原子1個は、どれくらい小さい？}{\bklevelrange{1}{2}}{スマホと原子炉、どちらも原子の物理で動く？}
\bkmichishirube{2-7,2-10,3-12,熱-16,波-13,電-2}{原-1,原-2,原-10,原-11,原-15}

\begin{goal}
\begin{itemize}[label={\color{bkTheme}●}, leftmargin=1.6\zw, itemsep=0pt, topsep=1pt]
\item \bklevel{1}\enspace 原子1個の大きさと、1円玉の中の原子の数を、けたで言える。スマホのチップと原子炉が、原子のどの部分の物理で動くかを説明できる。
\item \bklevel{2}\enspace 何を無視したかを書いてから、原子の数・1個の大きさ・1個あたりのエネルギーを、7つの段で求められる。
\item \bklevel{1}\enspace 身近な「なぜ？」の答えが、この巻のどの章にあるかを探せる。
\end{itemize}
\end{goal}

\begin{hatena}
\noindent\begin{minipage}[c]{0.56\linewidth}\raggedright
1円玉は、アルミニウムだけでできていて、重さはちょうど\qty{1}{g}です。この1円玉を、半分に、また半分にと、どこまでも細かく分けていくと、最後には、それ以上分けるとアルミニウムでなくなる小さな粒、\textbf{原子}にたどりつきます。\par\vspace{3pt}
原子1個は、どれくらい小さいのでしょう。1円玉1枚の中に、原子はいくつあるのでしょう。
\end{minipage}\hfill
\begin{minipage}[c]{0.41\linewidth}\centering
\begin{tikzpicture}[x=0.86mm, y=0.86mm]
  \filldraw[fill=bkGray!22, draw=bkGray!80!black, line width=0.8pt] (12,14) circle (11);
  \draw[bkGray!70!black, line width=0.4pt] (12,14) circle (9.6);
  \node[font=\sffamily\bfseries\large, text=bkGray!80!black] at (12,14) {1};
  \node[txa lab, anchor=north] at (12,2) {1円玉（\qty{1}{g}）};
  \draw[bkGray!80!black, line width=0.5pt, dashed] (17,18) -- (36,27) (17,10) -- (36,1);
  \filldraw[fill=white, draw=bkTheme, line width=0.9pt] (36,1) rectangle (54,27);
  \foreach \i in {0,...,3} {\foreach \j in {0,...,4} {\txaatom{({38.8+\i*4.15},{3.8+\j*4.85})}{1.9}{bkGray}}}
  \node[font=\sffamily\bfseries\Large, text=phHatena!80!black, fill=white, inner sep=1pt] at (45,14) {?};
  \node[txa lab, anchor=south] at (45,27.6) {原子の粒};
\end{tikzpicture}
\end{minipage}
\end{hatena}

この本は、身の周りの「原子」の「なぜ？」に、物理で答える本です。スマホのチップは、ケイ素の原子の中の電子を、電圧でたくみに動かして計算しています。原子炉は、ウランの原子の中心にある原子核を2つに分けて、熱を取り出しています。どちらも、目に見えないほど小さな原子の物理で動いているのです。序章では、原子がどれくらい小さいか、原子から取り出すエネルギーがどれくらいの大きさかを、数で確かめながら、この本で何を学ぶか、どの章で学ぶかを見ていきます。詳しい計算は、それぞれの章でします。いまはわからない式やことばが出てきても、だいじょうぶです。

\begin{yattemiyou}[紙を半分に、また半分に切っていこう]
\textbf{用意するもの}：はば\qty{1}{cm}・長さ\qty{20}{cm}の紙の帯、はさみ、ものさし、電卓\par\vspace{2pt}
\noindent\begin{minipage}[c]{0.62\linewidth}
\begin{enumerate}[label={\bkcircled[bkGreen]{\arabic*}}, leftmargin=2.4\zw, itemsep=1pt, topsep=0pt]
\item 紙の帯を長さが半分になるように切り、片方を、また半分に切る。これを繰り返すと、何回まで切れるだろうか。切るたびに、長さをはかって書きとめる。
\item 電卓で $20\div2\div2\div\cdots$ と、2でわり続ける。何回わると\qty{1}{mm}（\qty{0.1}{cm}）より短くなるだろうか。
\item 原子1個の大きさを約\qty{0.26}{nm}とする（\qty{1}{cm}の1億分の2.6）。\qty{20}{cm}を何回わると、これより短くなるだろうか。
\end{enumerate}
\end{minipage}\hfill
\begin{minipage}[c]{0.35\linewidth}\centering
\begin{tikzpicture}[x=0.8mm, y=1mm]
  \foreach \k/\w in {0/40,1/20,2/10,3/5,4/2.5} {
    \filldraw[fill=bkYellow!35, draw=bkInk, line width=0.5pt] (0,-\k*6) rectangle (\w,-\k*6+3.2);}
  \node[txa lab, anchor=west] at (41,1.6) {\qty{20}{cm}};
  \node[txa lab, anchor=west] at (21,-4.4) {\qty{10}{cm}};
  \node[txa lab, anchor=west] at (11,-10.4) {\qty{5}{cm}};
  \node[txa lab, anchor=west] at (6,-16.4) {\qty{2.5}{cm}};
  \node[txa lab, anchor=west] at (3.5,-22.4) {……};
\end{tikzpicture}
\end{minipage}\par\vspace{5pt}
\noindent\textbf{約束}：はさみの刃を、人や自分の手に向けない。小さくなった紙の切れはしを、口に入れない。
\end{yattemiyou}

\noindent やってみようの答え：①はさみで切れるのは、ふつう7〜8回くらいまでです。7回で約\qty{1.6}{mm}、8回で約\qty{0.8}{mm}になり、指でつまめなくなります。②8回わると\qty{0.078}{cm}で、\qty{1}{mm}より短くなります。③29回わると約\qty{0.37}{nm}、30回わると約\qty{0.19}{nm}です。\textbf{約30回}、半分にすれば、原子1個の大きさになります。$2^{10}=1024$ なので、30回半分にすると、長さは約 $(1000)^3=10$億分の1になるのです（$2^{30}\approx1.07\times10^9$。2-7）。とちゅうの20回めで約\qty{0.19}{\micro m}になり、ここで、光の顕微鏡で見分けられる限界（約\qty{0.2}{\micro m}。波-13）より小さくなります。原子は、光でも見えないほど小さいのです。数学の本の2-7では、紙を42回折ると、月まで届くほど厚くなることを計算しました。倍々の数は、大きくするほうにも、小さくするほうにも、あっという間にけたが進みます。

\begin{omoidasou}
\bkrecall{2-7}{同じ数を何回もかけるのが累乗。$10^{-9}$ は10億分の1。体積は長さの3乗で、1辺を2倍にすると体積は $2^3=8$ 倍（2-7②）。}
\bkrecall{2-10}{単位の前の k（キロ）は1000倍、M（メガ）は100万倍、m（ミリ）は1000分の1、μ（マイクロ）は100万分の1、n（ナノ）は10億分の1（2-10③）。}
\bkrecall{3-12}{けたを数えるのが常用対数。$\log_{10}1000=3$。目もりを1つ進むごとに10倍にする目もりが、対数の目もり（3-12①⑥）。}
\bkrecall{熱-16}{炭素（石炭のおもな成分）\qty{1}{mol}が燃えると、二酸化炭素ができ、約\qty{393.5}{kJ}の熱が出る。}
\bkrecall{電-2}{電子を\qty{1}{V}の電圧で加速したときに増える運動エネルギーが、\qty{1}{eV}（電子ボルト）。$\qty{1}{eV}=\qty{1.60e-19}{J}$（電-2 ⑥）。}
\bkrecall{波-13}{目に見える光の波長は、約\qtyrange{380}{780}{nm}。光の顕微鏡で見分けられるのは、約\qty{0.2}{\micro m}まで。}
\end{omoidasou}

%% ================================================================
\Needspace{16\baselineskip}
\bksection{hon}{この巻の地図}

\bksubsection{1}{原子は、どれくらい小さい？}
\leavevmode\label{ph:原-0:sec1}原子の大きさを言うには、とても小さな長さの単位が要ります。この本では、ナノメートル（$\qty{1}{nm}=10^{-9}$\,m。\qty{1}{mm}の100万分の1）と、フェムトメートル（$\qty{1}{fm}=10^{-15}$\,m。\qty{1}{nm}の100万分の1）をよく使います。\bkidx{なのめーとる}{ナノメートル}図\ref{ph:原-0:fig-length}は、1円玉から原子核までを、1つのものさしに並べたものです。目もりは、1つ右へ進むごとに10倍になります（3-12 の対数の目もり）。
\begin{bkfigure}
\begin{tikzpicture}[x=1mm, y=1mm]
  % x = 9*(log10(L/m)+15)
  \draw[bkInk, line width=0.8pt, -{Stealth[length=2.4mm]}] (-2,0) -- (130,0);
  \foreach \k in {0,...,14} {\draw[bkInk, line width=0.4pt] (\k*9,0) -- (\k*9,-1.2);}
  \foreach \x/\t in {0/1\,fm, 27/1\,pm, 54/1\,nm, 81/1\,\textmu m, 108/1\,mm} {\draw[bkInk, line width=0.6pt] (\x,0) -- (\x,-2); \node[txa lab, anchor=north] at (\x,-2.4) {\t};}
  \node[txa lab, anchor=north] at (126,-2.4) {\qty{10}{cm}};
  \draw[bkInk, line width=0.6pt] (126,0) -- (126,-2);
  % 原子核（2〜15 fm）
  \draw[bkRed!70, line width=3pt] (2.7,2.2) -- (10.6,2.2);
  \node[txa lab, align=center, anchor=south] at (6.5,4) {原子核\\数 fm};
  % 原子（0.1〜0.3 nm）
  \draw[bkBlue!70, line width=3pt] (45,2.2) -- (49.3,2.2);
  \node[txa lab, align=center, anchor=south] at (47,4) {原子\\0.1〜0.3\,nm};
  % チップの細い所 10 nm
  \fill[bkPurple] (63,0) circle (1.3);
  \node[txa lab, align=right, anchor=north east] at (65,-8) {チップの細い所\\約 \qty{10}{nm}};
  % 顕微鏡の限界 0.2 μm と可視光
  \fill[bkGreen!60!black] (74.7,0) circle (1.3);
  \node[txa lab, align=left, anchor=north west] at (72.5,-8) {光の顕微鏡の限界\\約\qty{0.2}{\micro m}};
  \draw[bkOrange!90, line width=3pt] (77.2,2.2) -- (80,2.2);
  \node[txa lab, align=center, anchor=south] at (80,4) {目に見える\\光の波長};
  % 1円玉 2 cm
  \fill[bkGray!80!black] (119.7,0) circle (1.3);
  \node[txa lab, align=center, anchor=south] at (117,3.5) {1円玉の直径\\\qty{2}{cm}};
\end{tikzpicture}
\bkfigcaption{長さのものさし。目もりが1つ進むごとに10倍。丸は1つの値、太い線は範囲。原子は、光の波長の千分の1ほど。原子核は、原子のさらに数万分の1。}\label{ph:原-0:fig-length}
\end{bkfigure}
\noindent 原子は、光の顕微鏡で見える限界のさらに約1000分の1です。そのまん中にある\textbf{原子核}は、原子のさらに数万分の1しかありません（金の原子で約2万分の1。原-1 ⑤）。\bkidx{げんしかく}{原子核}原子の大きさは、目で見ずに、どうやってわかるのでしょう。じつは、重さと密度と、\textbf{アボガドロ定数}（\qty{1}{mol}の中の粒の数 $N_{\mathrm A}=6.02\times10^{23}$）がわかれば、割り算だけで出せます（原-1 ⑦）。1円玉で確かめてみましょう。

\phboxhead{例題の解き方：7つの段}
この本の例題は、いつも同じ7つの段で解きます。①場面、②何を無視するか（理想化）、③図と記号、④使う式を選ぶ理由、⑤式を立てる、⑥解く、⑦答えを確かめる（単位・大きさ・極端な場合）です。数学の本の例題よりも段が多いのは、式を立てる前の「モデルを作る」部分と、答えを出したあとの「確かめる」部分を、ていねいに書くためです。何を無視したかは、式を立てる前に、次のような「モデル化の札」にも書きます。
\begin{phmodel}[原子を、すき間なく積んだ立方体と考える]
\phignore{原子の丸い形と、原子どうしの間のすき間（すき間も、その原子の場所にふくめる）。1円玉の表面の模様。ほかの原子のまじり。}
\phassume{1円玉は、アルミニウムの原子だけが、どこでも同じ混み具合で集まってできている。}
\phlimit{原子の「大きさ」は、はかり方で少し変わる。ここで出すのは、1個が占める場所の大きさ。}
\end{phmodel}
\bkidx{もでる}{モデル}このように、知りたいことに関係の小さいものを思い切って無視して作った、簡単にした考え方を\textbf{モデル}といいます。

\begin{phreidai}[1円玉の中の原子の数と大きさ（\bkhatenaref）]\label{ph:原-0:ex1}
1円玉は、アルミニウムだけでできていて、質量は\qty{1.00}{g}です。アルミニウム\qty{1}{mol}の質量は\qty{27.0}{g}、密度は\qty{2.70}{g/cm^3}です。1円玉の中の原子の数と、原子1個が占める場所を立方体とみたときの1辺を求めましょう。アボガドロ定数は $N_{\mathrm A}=6.02\times10^{23}$\,/mol とします。
\phbamen
目に見える1円玉の重さと体積から、目に見えない原子1個の大きさを求めます。
\phmushi
モデル化の札のとおり、原子の丸い形とすき間を無視し、1個が占める場所を、すき間なく積んだ立方体とします。
\phzu
質量 $m=\qty{1.00}{g}$、モル質量 $M=\qty{27.0}{g/mol}$、密度 $\rho=\qty{2.70}{g/cm^3}$、原子の数 $N$、1個の体積 $v$、1辺 $a$。
\phriyu
\qty{1}{mol}の中に $N_{\mathrm A}$ 個あるので、何 mol かがわかれば数が出ます。体積は質量÷密度で、それを $N$ 個で等分すれば1個の体積です。立方体の体積は1辺の3乗なので、1辺は3乗根で求めます（2-7②・3-11⑤）。
\phtateru
\[N=\frac{m}{M}\,N_{\mathrm A},\qquad V=\frac{m}{\rho},\qquad v=\frac{V}{N},\qquad a=\sqrt[3]{v}\]
\phtoku
\[N=\frac{1.00}{27.0}\times6.02\times10^{23}\approx2.23\times10^{22}\approx\qty{2.2e22}{個}\]
\[V=\frac{1.00}{2.70}\approx\qty{0.370}{cm^3}=\qty{3.70e-7}{m^3}\]
\[v=\frac{3.70\times10^{-7}}{2.23\times10^{22}}\approx\qty{1.66e-29}{m^3}\]
\[a=\sqrt[3]{1.66\times10^{-29}}\approx\qty{2.55e-10}{m}\approx\qty{0.26}{nm}\]
\phtashikame
単位：g÷(g/mol)×(1/mol)＝個（数）。$\mathrm{m^3}$ の3乗根は m。\\
大きさ：原-1 ⑦で水の分子1個は約\qty{0.31}{nm}でした。原子や小さな分子は、どれも\qtyrange{0.1}{0.3}{nm}くらいで、合っています。\\
極端な場合：1円玉を2枚にすると、原子の数は2倍ですが、1個の大きさは変わりません。大きさは、1 mol の質量と密度だけで決まります。
\kotae 約 $2.2\times10^{22}$ 個。1辺 約\qty{0.26}{nm}
\end{phreidai}
\noindent 例題の計算を、一般の形にまとめておきます。1個の体積は $v=V/N=(m/\rho)\big/(mN_{\mathrm A}/M)$ で、1円玉の質量 $m$ は消えます。
\begin{phflow}[物の重さから、原子1個の大きさへ]
\phflowrow{原子の数\\$N=\dfrac{m}{M}\,N_{\mathrm A}$}{1種類の原子だけ。\\すき間も、1個の\\場所にふくめる}{1個の体積\\$v=\dfrac{V}{N}=\dfrac{M}{\rho N_{\mathrm A}}$}
\phflowrow{$v=\dfrac{M}{\rho N_{\mathrm A}}$}{1個の場所を\\立方体とみなす}{1辺\\$a=\left(\dfrac{M}{\rho N_{\mathrm A}}\right)^{1/3}$}
\end{phflow}

\begin{phbaisuru}{原子1個の1辺 $a$}
もし、アルミニウムの原子の1辺が2倍だったら、1個が占める場所の体積は $2^3=8$ 倍になります（2-7②）。1円玉の体積は同じなので、入る原子の数は $\tfrac18$ の約 $2.8\times10^{21}$ 個になります。反対に、1辺が半分なら、数は8倍です。数は、1辺の3乗に反比例します。
\begin{bkcenter}
\begin{tikzpicture}[x=1mm, y=1mm]
  \txacube{(0,0)}{12}{bkBlue}
  \draw[bkBlue!70!black, line width=0.4pt] (6,0) -- (6,12) (0,6) -- (12,6) (6,12) -- (10.8,15.6) (2.4,13.8) -- (14.4,13.8) (12,6) -- (16.8,9.6) (14.4,1.8) -- (14.4,13.8);
  \node[txa lab, align=center, anchor=north] at (8,-2) {1辺 $a$ の原子が\\8個};
  \draw[bkInk, txa arr] (25,9) -- (36,9);
  \txacube{(44,0)}{12}{bkOrange}
  \node[txa lab, align=center, anchor=north] at (52,-2) {1辺 $2a$ の原子なら\\同じ場所に1個};
  \node[txa lab, anchor=west, text width=46mm, align=left] at (74,9) {1辺が2倍 → 体積は8倍\\→ 1円玉の中の数は $\tfrac18$};
\end{tikzpicture}
\bkfigcaption{1辺を2倍にすると、1個の体積は8倍。同じ場所に入る数は8分の1になる。}
\end{bkcenter}
\end{phbaisuru}

\Needspace{6\baselineskip}
\phboxhead{\bkhatenaref の答え}
\label{ph:原-0:hatena}1円玉1枚の中には、アルミニウムの原子が約 $2.2\times10^{22}$ 個あります。原子1個が占める場所は、1辺 約\qty{0.26}{nm}（$2.6\times10^{-10}$\,m）の立方体ほどです。\qty{1}{mm}の中に、原子が約400万個並ぶ大きさです。

どれほどの数かを、長さで感じてみましょう。1円玉の原子を、1個\qty{0.26}{nm}ずつ1列に並べると、
\[2.2\times10^{22}\times0.26\times10^{-9}\ \mathrm{m}\approx5.7\times10^{12}\ \mathrm{m}=5.7\times10^{9}\ \mathrm{km}\]
で、約57億 km になります。地球から月までの距離（約38万 km）の約1万5千倍です。光の速さで進んでも、5時間以上かかります。1個1個はとても小さく、数はとても多い。これが原子です。

%% ================================================================
\Needspace{14\baselineskip}
\bksubsection{2}{エネルギーのけた：燃える と 分かれる}
\leavevmode\label{ph:原-0:sec2}原子1個のエネルギーを、J（ジュール）で書くと、$10^{-19}$ のような小さな数ばかりになります。そこで、この本では\textbf{電子ボルト}（eV）という単位をよく使います。\bkidx{でんしぼると}{電子ボルト}電子を\qty{1}{V}の電圧で加速したときに増える運動エネルギーが、\qty{1}{eV}です（電-2 ⑥）。
\pheq{\qty{1}{eV}=\qty{1.602176634e-19}{J}\approx\qty{1.60e-19}{J}}{eV と J の換算}
\noindent 原子の中の電子がかかわるエネルギーは、数 eV ほどです。原子核がかかわるエネルギーは、その100万倍の MeV（メガ電子ボルト。$\qty{1}{MeV}=10^6$\,eV）のけたになります。図\ref{ph:原-0:fig-energy}は、この本に出てくるエネルギーを、1つのものさしに並べたものです。
\begin{bkfigure}
\begin{tikzpicture}[x=1mm, y=1mm]
  % x = 11*(log10(E/eV)+2)
  \draw[bkInk, line width=0.8pt, -{Stealth[length=2.4mm]}] (-2,0) -- (124,0);
  \foreach \k in {0,...,11} {\draw[bkInk, line width=0.4pt] (\k*11,0) -- (\k*11,-1.2);}
  \foreach \x/\t in {0/0.01\,eV, 22/1\,eV, 44/100\,eV, 66/10\,keV, 88/1\,MeV, 110/100\,MeV} {\draw[bkInk, line width=0.6pt] (\x,0) -- (\x,-2); \node[txa lab, anchor=north] at (\x,-2.4) {\t};}
  \fill[bkGray!80!black] (4.5,0) circle (1.3);
  \node[txa lab, align=center, anchor=south] at (4.5,3.5) {室温の\\熱 $k_{\mathrm B}T$\\\qty{0.026}{eV}};
  \draw[bkOrange!90, line width=3pt] (24.2,2.2) -- (27.7,2.2);
  \node[txa lab, align=center, anchor=south] at (21,4.2) {目に見える\\光の光子\\1.6〜3.3\,eV};
  \fill[bkRed!80!black] (28.7,0) circle (1.3);
  \node[txa lab, align=center, anchor=north] at (28.7,-8) {炭素1個が\\燃える\\\qty{4.1}{eV}};
  \fill[bkBlue!80!black] (34.5,0) circle (1.3);
  \node[txa lab, align=left, anchor=south west] at (32,3.5) {水素の\\電離\\\qty{13.6}{eV}};
  \fill[bkPurple] (74.6,0) circle (1.3);
  \node[txa lab, align=center, anchor=south] at (74.6,3.5) {X線写真の\\X線\\\qty{60}{keV}};
  \fill[bkGreen!60!black] (86,0) circle (1.3);
  \node[txa lab, align=center, anchor=north] at (86,-8) {γ線\\（セシウム137）\\\qty{0.662}{MeV}};
  \fill[bkRed!80!black] (113.3,0) circle (1.3);
  \node[txa lab, align=center, anchor=south] at (113.3,3.5) {ウラン235\\1個の分裂\\\qty{200}{MeV}};
  \draw[bkBlue!70!black, line width=0.6pt] (2,-25) -- (2,-27) -- (40,-27) -- (40,-25);
  \node[txa lab, anchor=north, text=bkBlue!70!black] at (21,-27.5) {電子がかかわる（eV の世界）};
  \draw[bkRed!70!black, line width=0.6pt] (84,-25) -- (84,-27) -- (118,-27) -- (118,-25);
  \node[txa lab, anchor=north, text=bkRed!70!black] at (101,-27.5) {原子核がかかわる（MeV の世界）};
\end{tikzpicture}
\bkfigcaption{エネルギーのものさし（1個の粒・1回の反応あたり）。目もりが1つ進むごとに10倍。光子・電離は原-2・原-4、X線とγ線は原-13、核分裂は原-15 で学ぶ。}\label{ph:原-0:fig-energy}
\end{bkfigure}

\begin{phreidai}[炭素が燃える と ウランが分かれる]\label{ph:原-0:ex2}
炭素\qty{1}{mol}が燃えると\qty{393.5}{kJ}の熱が出ます。ウラン235 の原子核1個が分裂すると、約\qty{200}{MeV}のエネルギーが出ます。(1)\enspace 炭素の原子1個が燃えるときのエネルギーは何 eV ですか。分裂1回はその何倍ですか。(2)\enspace \qty{1}{kg}あたりでは、それぞれ何 J ですか。炭素とウラン235 の\qty{1}{mol}の質量は\qty{12.0}{g}・\qty{235}{g}、$N_{\mathrm A}=6.02\times10^{23}$\,/mol、$\qty{1}{eV}=\qty{1.60e-19}{J}$ とします。
\phbamen
火力発電所は炭素を燃やし、原子力発電所はウランを分裂させて、どちらも熱で水蒸気をつくります（熱-16・原-15）。1個の原子から出るエネルギーをくらべます。
\phmushi
炭素は全部が燃え、ウラン235 は全部が分裂するとします。燃料のほかの成分（石炭の中の水素や灰、燃料の中のウラン238）は無視します。
\phzu
1個あたりのエネルギー $E_{\text{燃}}$・$E_{\text{分}}$〔eV〕、\qty{1}{kg}あたりのエネルギー〔J〕。
\phriyu
\qty{1}{mol}のエネルギーを $N_{\mathrm A}$ でわれば、1個あたりになります。\qty{1}{kg}あたりは、「\qty{1}{kg}の中の原子の数×1個あたり」です。
\phtateru
\[E_{\text{燃}}=\frac{\qty{393.5e3}{J}}{N_{\mathrm A}},\qquad \text{1 kg の原子の数}=\frac{1000}{M}\,N_{\mathrm A}\]
\phtoku
(1)
\[E_{\text{燃}}=\frac{3.935\times10^5}{6.02\times10^{23}}\approx\qty{6.54e-19}{J},\qquad \frac{6.54\times10^{-19}}{1.60\times10^{-19}}\approx\qty{4.1}{eV}\]
\[\frac{E_{\text{分}}}{E_{\text{燃}}}=\frac{2.00\times10^8}{4.1}\approx4.9\times10^7\ \text{倍}\]
(2)\enspace 炭素：
\[\frac{1000}{12.0}\times\qty{393.5e3}{J}\approx\qty{3.3e7}{J}\]
ウラン235：
\[\frac{1000}{235}\times6.02\times10^{23}\times200\times10^6\times1.60\times10^{-19}\approx\qty{8.2e13}{J}\]
比は $8.2\times10^{13}\div(3.3\times10^7)\approx2.5\times10^6$ 倍です。
\phtashikame
単位：(J/mol)÷(1/mol)＝J。MeV は $10^6$ eV。\\
大きさ：1回あたりは約5千万倍ですが、\qty{1}{kg}あたりは約250万倍です。ウランの原子1個は炭素の約20倍重く、\qty{1}{kg}の中の原子の数が約20分の1だからです（$4.9\times10^7\div19.6\approx2.5\times10^6$）。原-15 例題1 の「ウラン235 の\qty{1}{g}で $\qty{8.2e10}{J}$」と合います。\\
極端な場合：エネルギーの大きさは、原子の数ではなく、1回の反応で何が変わるかで決まります。電子の結びつきが変わる化学の反応は eV、原子核の中の結びつきが変わる核分裂は MeV です。
\kotae (1)\enspace 約\qty{4.1}{eV}、約5千万倍。(2)\enspace 炭素 約 $\qty{3.3e7}{J}$、ウラン235 約 $\qty{8.2e13}{J}$（約250万倍）
\end{phreidai}
\noindent 原子核の中の結びつきがなぜそれほど強いのか、分裂でなぜエネルギーが出るのかは、$E=mc^2$ と結合エネルギー（原-11 ②④⑤）で学びます。原子炉で分裂の数をどうやって一定に保つかは、原-15 で学びます。

%% ================================================================
\Needspace{8\baselineskip}
\bksubsection{3}{スマホのチップと原子炉：原子の2つの出口}
\leavevmode\label{ph:原-0:sec3}原子は、まん中の小さな原子核と、その周りの電子でできています（原-1）。この本で学ぶ原子の物理には、大きく2つの出口があります（図\ref{ph:原-0:fig-exit}）。
\begin{itemize}[leftmargin=1.2\zw, itemsep=2pt, topsep=2pt, label={\color{bkTheme}●}]
\item \textbf{電子の出口}　原子の中の電子が、とびとびのエネルギーの「席」を移ると、光が出たり吸われたりします（原-4）。原子どうしが電子を出し合うと、化学の結びつきができます（原-7）。固体の中では、電子の席が帯（バンド）になり、それで金属・半導体・絶縁体が分かれます（原-8）。ケイ素の半導体で作ったダイオードやトランジスタが、LED・太陽電池・スマホのチップです（原-9・原-10）。
\item \textbf{原子核の出口}　原子核の中の結びつきは、電子の結びつきより100万倍ほど強く、変わるときに MeV のエネルギーが出ます（原-11）。ひとりでにこわれる原子核は、放射線を出します（原-12〜14）。ウランの原子核を分裂させるのが原子炉で、水素の原子核をくっつける（核融合）のが太陽です（原-11・原-15）。
\end{itemize}
\begin{bkfigure}
\begin{tikzpicture}[x=1mm, y=1mm]
  \fill[bkBlue!8] (64,0) circle (17);
  \draw[bkBlue!60!black, line width=0.6pt, dashed] (64,0) circle (17);
  \draw[bkBlue!60!black, line width=0.4pt, dashed] (64,0) circle (10);
  \foreach \a in {30,150,270} {\fill[bkBlue!80!black] ($(64,0)+(\a:17)$) circle (1.1);}
  \foreach \a in {90,210,330} {\fill[bkBlue!80!black] ($(64,0)+(\a:10)$) circle (1.1);}
  \fill[bkRed!80!black] (64,0) circle (1.6);
  \node[txa lab, anchor=north] at (64,-19) {原子（電子と原子核）};
  \node[txa lab, anchor=north, fill=bkBlue!8] at (64,-2.4) {原子核};
  \node[txa box, draw=bkBlue!70!black] (L) at (16,0) {\textbf{電子の出口}\\光・化学の結びつき\\半導体・チップ\\数 eV\\第1〜3部};
  \node[txa box, draw=bkRed!70!black] (R) at (112,0) {\textbf{原子核の出口}\\放射線・原子炉\\太陽の光\\数 MeV〜200 MeV\\第4部};
  \draw[bkBlue!70!black, txa arr] (46,0) -- (L.east);
  \draw[bkRed!70!black, txa arr] (82,0) -- (R.west);
\end{tikzpicture}
\bkfigcaption{原子の物理の2つの出口。電子の点は、見やすいよう大きく描いた（電子の大きさは、測れる限り0）。原子核も、本当は原子の数万分の1の大きさ。}\label{ph:原-0:fig-exit}
\end{bkfigure}
\begin{phmijika}[身の回りの物の中で働く、原子の物理]
スマホと原子炉のほかにも、身の回りには、原子の物理で動く物がたくさんあります。
\begin{bkcenter}
\small
\begin{tblr}{colspec={Q[l,m,wd=34mm] X[l,m] Q[l,m,wd=17mm]}, width=\linewidth, rowsep=0.6pt, colsep=3pt,
  hlines={0.4pt,bkLine}, row{1}={bg=bkTheme!12, font=\sffamily\bfseries\footnotesize}}
身の回りの物 & 働いている原子の物理 & 学ぶ章 \\
花火の色 & 原子の電子が、とびとびのエネルギーの席を移って光を出す & 原-4 \\
スマホのカメラ & 光子1個が、電子1個を動けるようにする & 原-2 ⑫ \\
USBメモリ & 閉じこめた電子が、薄い壁をすりぬける & 原-6 \\
病院の MRI & 水素の原子核（陽子）の、向きが2つのスピン & 原-7 ⑥ \\
LED 照明・太陽電池 & 半導体の電子が、帯のすき間をこえて光子を出す・吸う & 原-9 \\
GPS & 速く動く衛星の時計は、相対論のとおりにずれる & 原-11 ⑦ \\
大昔の木や骨の年代 & 炭素14 の原子核が、半減期で減っていく & 原-12 \\
X線写真・PET 検査 & 光子が物にぶつかって減る、陽電子と電子が光子になる & 原-13 \\
\end{tblr}
\end{bkcenter}
\end{phmijika}
スマホのチップの中では、原子をいくつ並べたほどの大きさで電子を動かしているのでしょう。
\begin{phreidai}[チップの細い線に、ケイ素の原子はいくつ並ぶ？]\label{ph:原-0:ex3}
スマホのチップの中の、いちばん細い所のはばを\qty{10}{nm}とします。ケイ素の結晶では、となりあう原子の中心の間は\qty{0.235}{nm}です。このはばに、ケイ素の原子は何個並びますか。
\phbamen
ケイ素の結晶の上に作った、チップのとても細い線の細さを、原子の数で数えます（原-10）。
\phmushi
結晶の中では、原子はジグザグに並んでいますが、ここでは、となりあう原子を\qty{0.235}{nm}おきに1列にまっすぐ並べたとします。
\phzu
はば $L=\qty{10}{nm}$、原子の中心の間 $d=\qty{0.235}{nm}$、並ぶ数 $N$。
\phriyu
同じ間かくで1列に並ぶので、「はば÷間かく」が並ぶ数です。
\phtateru
\[N=\frac{L}{d}\]
\phtoku
\[N=\frac{10}{0.235}\approx42.6\approx43\ \text{個}\]
\phtashikame
単位：nm÷nm で、単位のない数（個）。\\
大きさ：例題1と同じ方法（1 mol の質量\qty{28.1}{g}、密度\qty{2.33}{g/cm^3}）で、ケイ素1個の場所の立方体の1辺を出すと約\qty{0.27}{nm}で、\qty{10}{nm}には約37個です。数え方で少し違いますが、どちらも「約40個」です。原-10 ⑫の「いちばん細い所が原子数十個ほど」と合います。\\
極端な場合：はばを半分の\qty{5}{nm}にすると、約21個です。
\kotae 約43個（約40個）
\end{phreidai}
\noindent 原子は割れないので、線を細くするのにも限りがあります（原-10 のまだわかっていないこと）。それに、原子が40個ほどの細さになると、電子は、ふつうなら通れない薄い壁をすりぬけてしまいます（トンネル効果。原-6）。チップを作る人は、原子の物理と向き合いながら、トランジスタを小さくしてきました。スマホのチップ1つには、こうしたトランジスタが100億個をこえて入っています（原-10 ②）。

%% ================================================================
\Needspace{10\baselineskip}
\bksubsection{4}{4つの部と、16のなぜ？}
\leavevmode\label{ph:原-0:sec4}この巻は、序章と4つの部（16の章）でできています。部の題は、原子が見せる4つの顔です。
\begin{itemize}[leftmargin=1.2\zw, itemsep=2pt, topsep=2pt, label={\color{bkTheme}●}]
\item \textbf{第1部　飛び出す}：電子・光・原子核が「飛び出して」くる実験から、光はつぶでもあり、電子は波でもあることを知る。原子が出す光と、とびとびのエネルギー（原-1〜4）。
\item \textbf{第2部　すりぬける}：電子の波をシュレディンガー方程式で書く。電子が壁をすりぬけるトンネル効果、原子の中の電子の席と周期表（原-5〜7）。
\item \textbf{第3部　切りかえる}：固体の中の電子の帯、半導体の pn 接合と LED・太陽電池、電流を切りかえるトランジスタと電子回路（原-8〜10）。
\item \textbf{第4部　分かれる}：質量とエネルギー、原子核の崩壊と放射線、放射線のはかり方と身の守り方、原子炉、原子の力とのつきあい方（原-11〜16）。
\end{itemize}
\begin{bkcenter}
\small
\begin{tblr}{colspec={Q[l,m,wd=19mm] X[l,m] Q[l,m,wd=37mm]}, width=\linewidth, rowsep=1.1pt, colsep=3pt,
  hlines={0.4pt,bkLine}, row{1}={bg=bkTheme!12, font=\sffamily\bfseries\footnotesize},
  row{2,7}={bg=bkTheme!6, font=\sffamily\bfseries\footnotesize}}
章 & 身近な「なぜ？」 & 答えのある所 \\
\SetCell[c=3]{l}{\phfaceicon[5mm]{1}\enspace 第1部　飛び出す：光と電子の二重性} & & \\
\phchapref{原-1} & 原子の中は、ほとんどからっぽだって本当？ & ⑤ ラザフォードの散乱 \\
\phchapref{原-2} & 熱した鉄は、なぜ温度で光の色が変わる？ & ⑪ プランクの式 \\
\phchapref{原-3} & 電子顕微鏡は、なぜふつうの顕微鏡より小さなものが見える？ & ⑥ 電子顕微鏡 \\
\phchapref{原-4} & 花火の色は、なぜ材料によって変わる？ & ③ 元素の指紋 \\
\SetCell[c=3]{l}{\phfaceicon[5mm]{2}\enspace 第2部　すりぬける：量子力学の入口} & & \\
\phchapref{原-5} & 電子は、原子の中のどこにいる？ & ④ 波動関数と確率 \\
\phchapref{原-6} & USBメモリは、どうやって電子を閉じこめて記録する？ & ④ フラッシュメモリ \\
\phchapref{原-7} & ネオンは、なぜほかの物質とほとんど結びつかない？ & ⑨ 満席の殻 \\
\end{tblr}
\end{bkcenter}
\begin{bkcenter}
\small
\begin{tblr}{colspec={Q[l,m,wd=19mm] X[l,m] Q[l,m,wd=37mm]}, width=\linewidth, rowsep=1.1pt, colsep=3pt,
  hlines={0.4pt,bkLine}, row{1,5}={bg=bkTheme!6, font=\sffamily\bfseries\footnotesize}}
\SetCell[c=3]{l}{\phfaceicon[5mm]{3}\enspace 第3部　切りかえる：固体と半導体} & & \\
\phchapref{原-8} & 金属は電気を通すのに、ガラスはなぜ通さない？ & ⑨ 帯のうまり方 \\
\phchapref{原-9} & LEDは、どうやって電気を直接、光に変える？ & ⑨ LED \\
\phchapref{原-10} & スマホのチップは、どうやって計算している？ & ⑦ CMOS の論理ゲート \\
\SetCell[c=3]{l}{\phfaceicon[5mm]{4}\enspace 第4部　分かれる：原子核と放射線} & & \\
\phchapref{原-11} & 太陽は、何をエネルギーのもとにして光っている？ & ⑪ 太陽の核融合 \\
\phchapref{原-12} & 大昔の木や骨の年代は、どうやってわかる？ & ⑦ 年代測定 \\
\phchapref{原-13} & X線写真で、なぜ骨が白くうつる？ & ③ 減弱と半価層 \\
\phchapref{原-14} & 目に見えない放射線を、どうやってはかる？ & ② 検出器 \\
\phchapref{原-15} & 原子炉は、どうやって反応を一定に保つ？ & ⑨ 遅れる中性子 \\
\phchapref{原-16} & 使い終わった核燃料は、どうする？ & ②③ 使い終わった燃料 \\
\end{tblr}
\bktabcaption{この巻の地図。「答えのある所」の丸数字は、その章の「物理の道具」の節の番号。}\label{ph:原-0:tab-chizu}
\end{bkcenter}
表\ref{ph:原-0:tab-chizu}は、それぞれの章で答える身近な「なぜ？」と、その答えがある所です。章の札の横の数は、その章が始まるページです。どの章から読んでもよいのですが、前の章の道具を使うので、はじめは順に読むのがおすすめです。とくに、第3部の半導体は第1部・第2部の電子の席を、第4部の原子炉は原-11〜13 の原子核と放射線を土台にしています。序章の\bkhatenaref の答えは、①にあります。身近な疑問から、数学と物理のどの章を順に読めばよいかを示した「なぜ？の道すじ」も、巻末にあります。

この巻は、物理シリーズの最後の巻です。『力学』『熱学』『波動』『電磁気学』で「原子の物理で」と先へ送った話（光のつぶ、とびとびのエネルギー、半導体、相対論など）は、この巻のそれぞれの章で答えます。そのあとの続きは、各章の「この先の学び」で、専門の本につなぎます。

この巻の物理の多くは、19世紀の終わりからの約50年の間に見つかりました（表\ref{ph:原-0:tab-nenpyo}）。目に見えない原子の中を、実験で1つずつ確かめていった歴史です。発見が技術になるまでも速く、核分裂の発見（1938年）から、はじめて連鎖反応を続けた原子炉（1942年）までは、わずか4年でした（原-15）。
\begin{bkcenter}
\small
\begin{tblr}{colspec={Q[c,m,wd=12mm] X[l,m] Q[l,m,wd=22mm]}, width=\linewidth, rowsep=0.6pt, colsep=3pt,
  hlines={0.4pt,bkLine}, row{1}={bg=bkTheme!12, font=\sffamily\bfseries\footnotesize}}
年 & できごと & 学ぶ章 \\
1895 & レントゲンが X線を見つける & 原-13 \\
1896 & ベクレルが、ウランから出る放射線を見つける & 原-12 \\
1897 & トムソンが、陰極線の正体（電子）の比電荷をはかる & 原-1 \\
1900 & プランクが、光のエネルギーのとびとびの粒（量子）を考える & 原-2 \\
1905 & アインシュタインが、光子と相対性理論を発表する & 原-2・原-11 \\
1911 & ラザフォードが、原子核を見つける & 原-1 \\
1913 & ボーアが、とびとびのエネルギーで水素の光を説明する & 原-4 \\
1924 & ド・ブロイが、電子も波だと考える & 原-3 \\
1926 & シュレディンガーが、電子の波の方程式を発表する & 原-5 \\
1932 & チャドウィックが、中性子を見つける & 原-1・原-11 \\
1938 & ハーンとシュトラスマンが、ウランの核分裂を見つける & 原-15 \\
1942 & フェルミの組が、はじめて連鎖反応を制御して続ける & 原-15 \\
1947 & ベル研究所で、トランジスタが動く & 原-10 \\
\end{tblr}
\bktabcaption{原子の物理の年表（この巻の章に書いた年）。}\label{ph:原-0:tab-nenpyo}
\end{bkcenter}

%% ================================================================
\Needspace{8\baselineskip}
\bksubsection{5}{数学の本とのつながり}
\leavevmode\label{ph:原-0:sec5}原子の物理の式の多くは、数学の本でもう計算しています。たとえば、半分ずつへる量（半減期）は 3-11 と 6-4、とびとびの値だけが許される問題（固有値）は 7-8 で扱いました。この巻では、同じ計算を繰り返すかわりに、「（6-4⑤）」のように数学の章と節を指します。そして、物理の本ならではの、次の3つを中心に書きます。
\begin{itemize}[leftmargin=1.2\zw, itemsep=2pt, topsep=2pt, label={\color{bkTheme}●}]
\item その式は、どんな実験と法則から出てくるのか。何を無視して（どんな仮定で）作ったのか。
\item 本物の物（原子・LED・チップ・放射線の測定器・原子炉）の値を入れると、何が起きるのか。
\item 仮定がくずれると、何が変わるのか。
\end{itemize}
\begin{bkcenter}
\small
\begin{tblr}{colspec={Q[l,m,wd=28mm] X[l,m] Q[l,m,wd=28mm]}, width=\linewidth, rowsep=1.1pt, colsep=3pt,
  hlines={0.4pt,bkLine}, row{1}={bg=bkTheme!12, font=\sffamily\bfseries\footnotesize}}
数学の本の章 & そこで解いたこと & この巻の章 \\
2-7・3-12 & 10の累乗、けたを数える対数 & 全部の章 \\
9-3・9-4 & 回る矢印で書く波、波の重ね合わせとスペクトル & 原-3・原-5・原-8 \\
7-8 & 固有値：とびとびの値だけが許される & 原-5・原-7・原-8 \\
5-14・6-11 & ふしぎな数 $e$、双曲線関数 & 原-6・原-9 \\
8-12 & ラプラシアン（球の形のとき） & 原-7・原-15 \\
3-12・2-17 & 対数のけた、2進数 & 原-9・原-10 \\
3-11・6-4 & 半分ずつへる量、時定数と半減期 & 原-12〜原-14 \\
6-14・9-9 & 連立微分方程式（親と子）、極と安定性 & 原-12・原-15 \\
3-14・3-17 & ばらつき：標準偏差、確率 & 原-12・原-14 \\
\end{tblr}
\bktabcaption{数学の本で解いた計算と、それを使うこの巻の章。}\label{ph:原-0:tab-math}
\end{bkcenter}
\noindent 数学の本をまだ読んでいなくても、だいじょうぶです。★が5つの内容をふくむ章は、「目でわかる」「手でわかる」「式でわかる」の3段に分けてあり、はじめの2段は、図と表計算を中心に、★3つまでの道具でたどれるように書いてあります。

%% ================================================================
\Needspace{8\baselineskip}
\bksubsection{6}{この本の読み方}
\leavevmode\label{ph:原-0:sec6}
\begin{bkfigure}
\begin{tikzpicture}[x=1mm, y=1mm]
  \txatalk{0}{sora}{ソラ}{原子って、1円玉の中に $10^{22}$ 個もあるんだね。そんなに小さい物を、どうやって調べたんだろう。}
  \txatalk{-20}{mina}{ミナ}{わたしは、光がつぶで、電子が波だって聞いて、ふしぎ……。目に見える物と、ふるまいが違うの？}
  \txatalk{-44}{delta}{デルタ}{原子は見えないから、飛んでくる電子や光を使って、中をのぞいたんだ（原-1）。小さな世界では、つぶと波の両方の顔が見えてくる（原-2・原-3）。★が5つの所は3段に分けてあるから、「ここまでで十分」の旗まで読めば、その章のゴールだよ。}
\end{tikzpicture}
\end{bkfigure}
\begin{taisetsu}
\def\tsl#1{\par\noindent\hangindent=6.2\zw\hangafter=1\makebox[6.2\zw][l]{\sffamily\bfseries\small\color{bkDeep}#1}}
\tsl{長さ}\qty{1}{nm}＝$10^{-9}$\,m、\qty{1}{fm}＝$10^{-15}$\,m。原子は約\qty{0.1}{nm}、原子核は数 fm。
\tsl{エネルギー}$\qty{1}{eV}=\qty{1.60e-19}{J}$、$\qty{1}{MeV}=10^6$\,eV。原子1個のエネルギーは、ふつう eV で書く。
\tsl{定数}$h=\qty{6.63e-34}{J.s}$、$c=\qty{3.00e8}{m/s}$、$N_{\mathrm A}=6.02\times10^{23}$\,/mol、$hc=1240$\,eV·nm、\qty{300}{K}の $k_{\mathrm B}T=\qty{0.0259}{eV}$。
\tsl{光の振動数}光の振動数は $f$ と書き、光子のエネルギーは $E=hf$（高校の教科書の $h\nu$ と同じ。原-2）。
\tsl{答えのけた}答えは、ふつう2けた。本物の測定の値とくらべる例題などは3〜4けたで書く。
\end{taisetsu}
\begin{itemize}[leftmargin=1.2\zw, itemsep=2pt, topsep=2pt, label={\color{bkTheme}●}]
\item \textbf{★の数}：内容の高さの目安です。★1つは「はじめの一歩」、★2つは「ふもとを歩く」、★3つは「山を登る」、★4つは「尾根を歩く」、★5つは「頂上をめざす」。物理の本の★は、物理の内容の高さで付けてあり、数学の本の★とは別の物差しです。
\item \textbf{章の札}：「7-8」のような札は数学の本の章、「原-4」のような札はこの本の章、「力-4」「熱-16」「波-13」「電-2」は第1巻『力学』・第2巻『熱学』・第3巻『波動』・第4巻『電磁気学』の章です。札の色は、その本の色です。本文で「（6-4⑤）」と書いたら、数学の本の6-4の⑤の節を見てください。
\item \textbf{物理の本の囲み}：「この式の読み方」「身近な例では」「〇〇を2倍にすると？」「まだわかっていないこと」「この先の学び」「一般の式から、よく使う式へ」などがあります。見本は、前付の「この本の使い方」にあります。
\item \textbf{本物の値}：定数、原子の質量、核種の半減期、半導体の値など、この巻で使う値は、巻末の「核種のデータ・物性値・定数」にまとめてあります。
\end{itemize}
\noindent 記号の読み方と意味は、巻末の「記号カード」にまとめてあります。同じ文字を、章によって別の意味に使うこともあります（たとえば $\lambda$ は、第1部では波長、第4部では崩壊定数）。そのときは、章のはじめに書いてあります。

%% ================================================================
\Needspace{8\baselineskip}
\bksubsection{7}{安全の約束}
\leavevmode\label{ph:原-0:sec7}この巻には、放射線・X線・レーザー・高い電圧・液体窒素など、正しく扱わないと体をきずつけるものが出てきます。そこで、この巻の「やってみよう」は、次の約束で作りました。
\begin{column}{安全：この本の「やってみよう」の約束}
\begin{enumerate}[label={\bkcircled{\arabic*}}, leftmargin=2.4\zw, itemsep=0pt, topsep=2pt]
\item \textbf{使うのは}：紙とはさみ、分光シート、乾電池（\qty{1.5}{V}を2本まで）と抵抗と LED、光電池、サイコロ、ドミノ、スマホの照度計アプリ、電卓・表計算・Python、自治体が公開している放射線の測定のデータ。
\item \textbf{見るだけにするもの}：放射線を出す物や装置、X線、高い電圧の真空の管、レーザー（ポインターも。のぞかない・人に向けない）、紫外線のランプ、液体窒素・ドライアイス、火。これらは、先生や科学館・病院の人がいる所で見るか、本で読むだけにします。
\item \textbf{放射線をはかるとき}：学校・自治体・科学館の測定器で、先生や係の人といっしょにはかります。放射線を出すかもしれない物（ある種の鉱物、古い夜光塗料の時計の文字盤など）を、自分で集めたり、こわしたりしません（原-11・原-14）。
\end{enumerate}
分光シートで照明を見るときも、太陽を直接見てはいけません。放射線の影響は「あるか、ないか」ではなく、受けた\textbf{量}で決まります。量のはかり方と、量を減らす方法は、原-14 で学びます。
\end{column}
\noindent 安全のための囲みは、各部のはじめの章（原-1・原-5・原-8・原-11）と、放射線をはかる原-14 にもあります。

%% ================================================================
\Needspace{6\baselineskip}
\bksubsection{8}{はてな？の答え：小さな原子と、2つの出口}
\leavevmode\label{ph:原-0:sec8}\bkhatenaref に、もう一度答えましょう。
\begin{enumerate}[label={\bkcircled[bkBlue]{\arabic*}}, leftmargin=2.4\zw, itemsep=2pt, topsep=2pt]
\item \textbf{原子の大きさ}　1円玉の原子は、1個が1辺 約\qty{0.26}{nm}の場所を占めます。質量・密度・アボガドロ定数から、割り算だけで出せました（①、\bkreidairef{1}）。
\item \textbf{原子の数}　1円玉1枚に、約 $2.2\times10^{22}$ 個。1列に並べると、月までの距離の約1万5千倍です（①）。
\item \textbf{原子核は、もっと小さい}　原子のまん中の原子核は、原子の数万分の1の大きさです。それでも、原子の質量のほとんどは原子核にあります（①、原-1 ⑤）。
\item \textbf{2つの出口}　電子がかかわるエネルギーは数 eV、原子核がかかわるエネルギーは MeV のけたです（②、\bkreidairef{2}）。スマホのチップは電子の物理、原子炉は原子核の物理で動きます（③）。
\end{enumerate}
つまり、\textbf{原子1個は1辺 約\qty{0.26}{nm}ほどで、1円玉1枚に約 $2.2\times10^{22}$ 個もある}ほど小さいのです。スマホのチップも、原子炉も、この小さな原子の中の物理で動いています。4つの部で、原子のそれぞれの顔を学びます（図\ref{ph:原-0:fig-kotae}）。
\begin{bkfigure}
\begin{tikzpicture}[x=1mm, y=1mm]
  \tikzset{tag/.style={draw=#1, line width=0.9pt, fill=white, rounded corners=1.6mm, text width=58mm, align=left, inner sep=1.8mm, font=\sffamily\footnotesize}}
  \node[tag=bkBlue, anchor=north west] (n1) at (0,0) {\phfaceicon[5mm]{1}\enspace\textbf{飛び出す}\\光はつぶ（光子）、電子は波でもある。原子は、決まった色の光を出す。\\\phchapref{原-2}\enspace\phchapref{原-4}};
  \node[tag=bkOrange!80!black, anchor=north west] (n2) at (66,0) {\phfaceicon[5mm]{2}\enspace\textbf{すりぬける}\\電子の波は、とびとびのエネルギーをとり、薄い壁をすりぬける。\\\phchapref{原-5}\enspace\phchapref{原-6}};
  \node[tag=bkGreen!70!black, anchor=north west] (n3) at (0,-27) {\phfaceicon[5mm]{3}\enspace\textbf{切りかえる}\\固体の電子の帯から、LED・太陽電池・チップのトランジスタへ。\\\phchapref{原-9}\enspace\phchapref{原-10}};
  \node[tag=bkPurple, anchor=north west] (n4) at (66,-27) {\phfaceicon[5mm]{4}\enspace\textbf{分かれる}\\原子核は、こわれたり分かれたりして、MeV のエネルギーと放射線を出す。\\\phchapref{原-11}\enspace\phchapref{原-15}};
\end{tikzpicture}
\bkfigcaption{原子の4つの顔。第1〜3部は電子の出口、第4部は原子核の出口。}\label{ph:原-0:fig-kotae}
\end{bkfigure}

\begin{phmada}[電子に、大きさはある？]
原子の大きさは約\qty{0.1}{nm}、原子核は数 fm とわかりました。では、原子の中の電子の大きさは、どれくらいでしょう。じつは、これまでのどんな実験でも、電子の大きさは見つかっていません（原-7 ⑥）。たとえば、電子がもつ小さな磁石の強さは、2023年に、1兆分の1ほどの細かさではかられました。大きさのない「点」として理論で計算した値と、そのけたまで合っています。電子は、本当に大きさのない点なのか、それとも、まだはかれないほど小さな大きさがあるのか。答えは、まだわかっていません。\par\vspace{2pt}
この巻の各章にも、「まだわかっていないこと」があります。陽子の大きさ（原-4）、「観測」とは何か（原-3）、高温超伝導のしくみ（原-8）、トランジスタはどこまで小さくできるか（原-10）、ニュートリノの質量（原-11）、\qty{100}{mSv}より少ない線量でがんは増えるか（原-14）、核融合で電気をつくれるか（原-15）などです。原子の物理は、いまも研究が続いています。
\end{phmada}

\begin{phketa}[原子の数とエネルギーの目安]
\phketaline{原子・原子核の大きさ}{\qtyrange{0.1}{0.3}{nm}・数 fm（原子の数万分の1）}
\phketaline{1円玉（\qty{1}{g}）の原子}{約 $2.2\times10^{22}$ 個、1個の場所の1辺 約\qty{0.26}{nm}（\bkreidairef{1}）}
\phketaline{室温の熱・目に見える光の光子}{\qty{0.0259}{eV}・\qtyrange{1.6}{3.3}{eV}}
\phketaline{炭素1個が燃える・水素の電離}{約\qty{4.1}{eV}・\qty{13.6}{eV}}
\phketaline{ウラン235 の1回の分裂}{約\qty{200}{MeV}（燃える炭素1個の約5千万倍。\bkreidairef{2}）}
\phketaline{チップの細い所（\qty{10}{nm}とした）}{ケイ素の原子 約40個（\bkreidairef{3}）}
\phboxhead{答えが桁で合っているか}
原子1個の大きさが nm より大きくなったり、1円玉の原子の数が $10^{23}$ をこえたりしたら、cm$^3$ と m$^3$、g と kg を取り違えています。1個の化学の反応で MeV、1個の核分裂で eV が出てきたら、J と eV の換算か、mol の扱いを取り違えています。
\end{phketa}

\begin{kotoba}\small
\bkword{原子}{物質をつくる小さな粒。大きさは約\qtyrange{0.1}{0.3}{nm}。まん中の原子核と、その周りの電子でできている。}
\bkword{原子核}{原子のまん中にある、とても小さく重い粒。陽子と中性子でできている。}
\bkword{アボガドロ定数}{\qty{1}{mol}の中にある粒の数。$N_{\mathrm A}=6.02\times10^{23}$\,/mol。}
\bkword{電子ボルト（eV）}{$\qty{1}{eV}=\qty{1.60e-19}{J}$。電子を\qty{1}{V}で加速したときに増える運動エネルギー。}
\bkword{電離}{原子から電子を引きはなして、イオンにすること。水素の原子では\qty{13.6}{eV}がいる。}
\bkword{核分裂}{重い原子核が、ほぼ半分の大きさの2つの原子核に分かれること。ウラン235 では1回に約\qty{200}{MeV}が出る。}
\bkword{モデル}{知りたいことに関係の小さいものを無視して作った、簡単にした考え方。}
\end{kotoba}

\begin{furikaeri}
\begin{checklist}
\item 原子1個の大きさと、1円玉の中の原子の数を、けたで言える。スマホのチップと原子炉が、原子のどの部分の物理で動くかを説明できる。\bkfurref{ph:原-0:sec3}
\item 何を無視したかを書いてから、原子の数・1個の大きさ・1個あたりのエネルギーを、7つの段で求められる。\bkfurref{ph:原-0:ex1}
\item 身近な「なぜ？」の答えが、この巻のどの章にあるかを探せる。\bkfurref{ph:原-0:sec4}
\end{checklist}
\end{furikaeri}
\bkfukidashi{sora}{原子は、1円玉に $10^{22}$ 個も入るほど小さいんだね。まずは原-1 で、見えない原子の中を、どうやってのぞいたかを調べよう。}
